\par\textbf{Desarrollo para comparar tu intento.} Al dividir la ley de Arrhenius a igual tensión, se cancela el prefactor y queda $\dot\varepsilon_2/\dot\varepsilon_1=\exp[-Q(1/T_2-1/T_1)/R]$. La diferencia de inversas es negativa al calentar, haciendo positivo el exponente.

\begin{enumerate}

\item Restar inversas de temperaturas absolutas.

\[\fitmath{\mathrm{inverseDifference}=\frac{1}{950}-\frac{1}{900}\approx -5.847953\times10^{-5}\,\mathrm{1/K}}\]

\item Convertir 250 kJ/mol a J/mol.

\[\fitmath{\mathrm{exponent}=\frac{-250000}{8.314462618}\cdot \left(\frac{1}{950}-\frac{1}{900}\right)\approx 1.758368\,}\]

\item Exponenciar para recuperar razón de tasas.

\[\fitmath{\mathrm{rateRatio}=\exp\left(\frac{-250000}{8.314462618}\cdot \left(\frac{1}{950}-\frac{1}{900}\right)\right)\approx 5.802958\,}\]

\end{enumerate}
